\begin{statement}[C.6.R4; Known with derivation: three exact leading factors]
For $i=3,5,7$, namely $j=287496,-3375,-884736$, one has
\[
 \operatorname{Sha}(E_i/\Q)=0,\qquad
 \frac{L'(E_i,1)}{\Omega_i}=4\widehat h_O(P_i).
\]
Thus the factors $q_3=q_5=q_7=4$ are exact, not numerical
rational reconstructions. For the smallest of these examples,
$E:y^2+y=x^3-38x+90$ and $P=(0,9)$, this gives the exact
reduction
\[
 N(P)=\exp\left(-\frac{L'(E,1)}{4\Omega_E}\right).
\]
\end{statement}
\begin{proof}
We use \citet[Theorem 1.2 and Table 9]{miller2011}. In the
notation of its Definition 1.1, the theorem proves
$\mathrm{BSD}(E/\Q,p)$ for every prime $p$ when the curve has
analytic rank at most one, conductor less than $5000$, and
its isogeny class is outside the eleven exceptions in Table 9.
The exception conductors are
\[
546,570,858,870,1050,1230,1938,1950,2370,2550,3270.
\]
The three conductors here are $3136,3136,361$, and their
analytic ranks are one by C.6.R2. Hence every hypothesis is
verified. The theorem is stated for all curves in the isogeny
class, using isogeny invariance, so no optimal-curve assumption
is being silently imposed on our minimal equations.

In all three rows
$2\prod_pc_p/\#E(\Q)_{\rm tors}^2=4$, and primitivity is
proved in C.6.R1. Therefore the quotient
$L'(E_i,1)/(4\Omega_i\widehat h_O(P_i))$ is the positive
integer $\#\operatorname{Sha}(E_i/\Q)$. To isolate this integer
without relying on an unbounded decimal approximation to a
canonical height, we give an elementary coarse height certificate.

After completing the square, write the duplication map on the
$x$-line as $[F(X,Z):G(X,Z)]$, with relatively prime integral
quartics. The program
\path{code/C_height_coarse_certificate.py} supplies integer
homogeneous cubics $A_0,B_0,A_1,B_1$ and $D\ne0$ with
\[
 A_0F+B_0G=DZ^7,\qquad A_1F+B_1G=DX^7.
\]
Let $C$ be the larger of the coefficient one-norms of $F,G$
and of the two pairs $(A_j,B_j)$. For coprime integers $a,b$,
the common divisor of $F(a,b),G(a,b)$ divides $D$, by these
identities. Taking the identity corresponding to the larger
of $|a|,|b|$ proves, together with the elementary upper bound,
\[
 |h(x(2Q))-4h(x(Q))|\leq\log C.
\]
Telescoping consequently gives
\[
 \left|\widehat h_O(P)-\frac{h(x(2^6P))}{2\cdot4^6}\right|
 \leq\frac{\log C}{6\cdot4^6}.
\]
The exact constants in row order $3,5,7$ are
\[
 C=26610617537591808,\quad188716953600,\quad285931133.
\]
Six rational doublings, the displayed identities, and directed
Arb logarithms give the quotient enclosures
$1\pm0.003$, $1\pm0.003$, and $1\pm0.009$, respectively,
after including the central-derivative error bound from C.6.R2.
Each interval contains only the integer one. All quartics,
Bezout coefficients, rational iterates and enclosures are
retained in the program's output and its JSON certificate.
Finally C.6.3 gives $A=1$ in row seven, proving the exponential
formula. Its algebraicity remains the separate question C.7(c).
\end{proof}
